\documentclass[11pt]{article}

\usepackage[margin=1in]{geometry}
\usepackage{amsmath,amssymb,amsthm,mathtools}
\usepackage{bm}

\title{Conditional Branched Autoregressive Q for Restaurant Actions}
\author{}
\date{}

\begin{document}
\maketitle

\section{Setup}
\label{sec:setup}

Let the structured restaurant action be
\[
a = (t, x, y, z),
\]
where
\begin{itemize}
  \item $t \in \mathcal{T}(s)$ is the action type,
  \item $x \in \mathcal{X}(s,t)$ is the first argument,
  \item $y \in \mathcal{Y}(s,t,x)$ is the second argument,
  \item $z \in \mathcal{Z}(s,t,x,y)$ is the third argument.
\end{itemize}

For the current restaurant environment, the mapping is
\[
t = \texttt{action\_type}, \qquad
x = \texttt{object1}, \qquad
y = \texttt{location}, \qquad
z = \texttt{object2}.
\]

The valid sets $\mathcal{T}(s)$, $\mathcal{X}(s,t)$, $\mathcal{Y}(s,t,x)$, and $\mathcal{Z}(s,t,x,y)$ are induced by the environment masks.

\section{Conditional Branched Joint Q}
\label{sec:joint-q}

We define a joint action-value function
\[
Q(s,a) = Q(s,t,x,y,z),
\]
and parameterize it autoregressively as
\begin{equation}
\label{eq:autoregressive-q}
Q(s,t,x,y,z)
= q_t(s,t)
+ q_x(s,t,x)
+ q_y(s,t,x,y)
+ q_z(s,t,x,y,z).
\end{equation}

This is a \emph{joint} Q-function because it scores the full structured action. It is also \emph{autoregressive} because each later term is conditioned on earlier choices. The current implementation in this repository uses the dueling-style decomposition described in Section~\ref{sec:dueling-form}.

\paragraph{Interpretation.}
\begin{itemize}
  \item $q_t(s,t)$ is the stage contribution for the action type: a learned scalar function of the state and the chosen action type (a joint score over $(s,t)$).
  \item $q_x(s,t,x)$ is a learned, advantage‑style contribution that jointly scores the state, action type, and first argument; it is a function of $(s,t,x)$ and should be read as a value (score), not as a conditional probability over $x$ given $(s,t)$.
  \item $q_y(s,t,x,y)$ is a learned contribution that jointly scores $(s,t,x,y)$ and captures the incremental value of choosing $y$ given the prefix.
  \item $q_z(s,t,x,y,z)$ is a learned contribution that jointly scores $(s,t,x,y,z)$ and captures the incremental value of choosing $z$ given the prefix.
\end{itemize}

Unlike an unconditional additive factorization, the argument terms here are joint score functions tied to the action prefix rather than statistical conditional distributions; they are not normalized over their arguments.

\section{Suffix Value Functions}
\label{sec:suffix-values}

Define the optimal suffix values recursively:
\begin{align}
V_x(s,t)
&=
\max_{x \in \mathcal{X}(s,t)}
\left[
q_x(s,t,x) + V_y(s,t,x)
\right], \\
V_y(s,t,x)
&=
\max_{y \in \mathcal{Y}(s,t,x)}
\left[
q_y(s,t,x,y) + V_z(s,t,x,y)
\right], \\
V_z(s,t,x,y)
&=
\max_{z \in \mathcal{Z}(s,t,x,y)}
q_z(s,t,x,y,z).
\end{align}

These are not separate environment values. They are the optimal continuation values induced by the autoregressive decomposition in Equation~\eqref{eq:autoregressive-q}.

\section{Greedy Decoding}
\label{sec:greedy-decoding}

The greedy action under this model is obtained by backward-consistent sequential maximization.

First choose the action type:
\begin{equation}
\label{eq:greedy-type}
t^*
=
\arg\max_{t \in \mathcal{T}(s)}
\left[
q_t(s,t) + V_x(s,t)
\right].
\end{equation}

Then choose the first argument:
\begin{equation}
\label{eq:greedy-x}
x^*
=
\arg\max_{x \in \mathcal{X}(s,t^*)}
\left[
q_x(s,t^*,x) + V_y(s,t^*,x)
\right].
\end{equation}

Then choose the second argument:
\begin{equation}
\label{eq:greedy-y}
y^*
=
\arg\max_{y \in \mathcal{Y}(s,t^*,x^*)}
\left[
q_y(s,t^*,x^*,y) + V_z(s,t^*,x^*,y)
\right].
\end{equation}

Finally choose the third argument:
\begin{equation}
\label{eq:greedy-z}
z^*
=
\arg\max_{z \in \mathcal{Z}(s,t^*,x^*,y^*)}
q_z(s,t^*,x^*,y^*,z).
\end{equation}

This yields the greedy action
\[
a^* = (t^*,x^*,y^*,z^*).
\]

\paragraph{Key property.}
The decoder is now mathematically aligned with the model:
\[
a^* \in \arg\max_a Q(s,a).
\]
That alignment is the main reason to prefer this formulation over unconditional branch-wise argmax.

\section{Bellman Target}
\label{sec:bellman-target}

For a transition $(s,a,r,s',d)$ with terminal flag $d \in \{0,1\}$, the standard Q-learning target is
\begin{equation}
\label{eq:bellman-target}
\hat{Q}(s,a)
=
r + \gamma (1-d) \max_{a'} Q_{\mathrm{tgt}}(s',a').
\end{equation}

Under the autoregressive factorization, the greedy next action is
\[
a'_* = (t'_*,x'_*,y'_*,z'_*),
\]
where each component is chosen by the same suffix-maximizing decoder as in Equations~\eqref{eq:greedy-type}--\eqref{eq:greedy-z}, but using $s'$ and the valid masks at $s'$.

Therefore
\begin{equation}
\label{eq:target-expanded}
\max_{a'} Q_{\mathrm{tgt}}(s',a')
=
Q_{\mathrm{tgt}}(s',t'_*,x'_*,y'_*,z'_*).
\end{equation}

For Double DQN, one would choose $a'_*$ using the online network and evaluate it using the target network:
\begin{equation}
\label{eq:double-dqn-target}
\hat{Q}(s,a)
=
r + \gamma (1-d) Q_{\mathrm{tgt}}(s',a'_*),
\qquad
a'_* \in \arg\max_{a'} Q_{\mathrm{online}}(s',a').
\end{equation}

\section{Dueling-Style Conditional Architecture}
\label{sec:dueling-form}

A dueling-style decomposition can be added in a way that remains compatible with the autoregressive joint Q semantics. The clean version is to use a single state-value baseline
\[
V(s),
\]
followed by centered conditional advantage terms at each stage:
\begin{align}
V(s) &= f_{\mathrm{value}}(h(s)), \\
A_t(s,t) &= g_t(h(s), t), \\
A_x(s,t,x) &= g_x(h(s), \phi_t(t), x), \\
A_y(s,t,x,y) &= g_y(h(s), \phi_t(t), \phi_x(x), y), \\
A_z(s,t,x,y,z) &= g_z(h(s), \phi_t(t), \phi_x(x), \phi_y(y), z).
\end{align}

The corresponding stage contributions are
\begin{align}
q_t(s,t)
&=
V(s)
+
\Bigl(
A_t(s,t) - \bar{A}_t(s)
\Bigr), \label{eq:q_t-dueling}\\
q_x(s,t,x)
&=
A_x(s,t,x) - \bar{A}_x(s,t), \label{eq:q_x-dueling}\\
q_y(s,t,x,y)
&=
A_y(s,t,x,y) - \bar{A}_y(s,t,x), \\
q_z(s,t,x,y,z)
&=
A_z(s,t,x,y,z) - \bar{A}_z(s,t,x,y),
\end{align}
where the centered baselines are computed over the valid masked choices at each stage:
\begin{align}
\bar{A}_t(s)
&=
\frac{1}{|\mathcal{T}(s)|}
\sum_{t' \in \mathcal{T}(s)} A_t(s,t'), \\
\bar{A}_x(s,t)
&=
\frac{1}{|\mathcal{X}(s,t)|}
\sum_{x' \in \mathcal{X}(s,t)} A_x(s,t,x'), \\
\bar{A}_y(s,t,x)
&=
\frac{1}{|\mathcal{Y}(s,t,x)|}
\sum_{y' \in \mathcal{Y}(s,t,x)} A_y(s,t,x,y'), \\
\bar{A}_z(s,t,x,y)
&=
\frac{1}{|\mathcal{Z}(s,t,x,y)|}
\sum_{z' \in \mathcal{Z}(s,t,x,y)} A_z(s,t,x,y,z').
\end{align}

The full joint Q remains
\[
Q(s,t,x,y,z)
= q_t(s,t) + q_x(s,t,x) + q_y(s,t,x,y) + q_z(s,t,x,y,z).
\]

\paragraph{Implementation.}
This formulation is implemented by the \texttt{RestaurantQNetwork} class in \texttt{anticipatory\_rl/agents/restaurant/dqn.py}.
The components are:
\begin{itemize}
    \item $h(s)$ from \texttt{encoder} (line 200).
    \item $V(s)$ from \texttt{value\_head} (line 206).
    \item $A_t(s,t)$ from \texttt{action\_type\_adv\_head} (line 207).
    \item $A_x, A_y, A_z$ from \texttt{object1\_adv\_head}, \texttt{location\_adv\_head}, and \texttt{object2\_adv\_head} respectively (lines 213--227).
    \item The centered contributions in Equations~\eqref{eq:q_t-dueling} and~\eqref{eq:q_x-dueling} are computed by the \texttt{*\_scores} methods (lines 232--284).
\end{itemize}

\paragraph{Why this is better than stage-specific value heads.}
It is tempting to introduce separate learned value terms such as $V_x(s,t)$ or $V_y(s,t,x)$ and then add them to the stage contributions. However, that creates an identifiability problem: if two learned terms have the same conditioning context and are simply added together, the model can shift mass between them without changing the total Q-value. Using a single global baseline $V(s)$ plus centered conditional advantages avoids that redundancy.

\paragraph{Interpretation.}
The terms $q_t, q_x, q_y, q_z$ should be understood as centered incremental contributions to the total joint Q-value, not as independent Bellman Q-functions. The actual continuation value after fixing a prefix is still given by the suffix maximization defined earlier in Section~\ref{sec:suffix-values}.

This can improve optimization, but it does not change the main semantic requirement: later stages must be conditioned on earlier ones, and decoding must remain aligned with the joint autoregressive Q.

\section{Restaurant-Specific Instantiation}
\label{sec:instantiation}

For the current environment, a natural specialization is
\begin{equation}
\label{eq:restaurant-q}
Q(s,\texttt{type},\texttt{obj1},\texttt{loc},\texttt{obj2})
=
q_{\mathrm{type}}(s,\texttt{type})
+ q_{\mathrm{obj1}}(s,\texttt{type},\texttt{obj1})
+ q_{\mathrm{loc}}(s,\texttt{type},\texttt{obj1},\texttt{loc})
+ q_{\mathrm{obj2}}(s,\texttt{type},\texttt{obj1},\texttt{loc},\texttt{obj2}).
\end{equation}

For action types that do not use some arguments, the corresponding valid set is a singleton containing the ``none'' token, so the same equations still apply without case-splitting.

\section{Why This Fixes the Prior Design}
\label{sec:why-fix}

The prior unconditional factorization effectively assumed that argument scores were reusable across unrelated action types. In contrast, the conditional branched model ensures:
\begin{itemize}
  \item object preferences are action-type dependent,
  \item location preferences are conditioned on the action type and earlier arguments,
  \item the greedy decoder is consistent with the modeled Q-function,
  \item the Bellman target is computed with respect to the same joint action semantics used at inference time.
\end{itemize}

\end{document}
